% year: תש"פ
% type: מבחן
% semester: א
% moed: ג
% part: ב
% number: 2
% points: 9
% categories: derivatives, mvt, multiple-choice
% source: exams/מבחנים/תשפ/מועד מיוחד תשפ.pdf

\begin{question}
תהי $f(x)$ פונקציה גזירה בקטע פתוח $(a,b)$ וידוע כי נגזרתה חיובית, כלומר $f'(x)>0$. מה מהבאים חייב להיות נכון?
\begin{enumerate}
  \item $f(x)$ פונקציה חסומה בקטע פתוח $(a,b)$.
  \item $f(x)$ פונקציה המקיימת $f(x)>0$ בקטע פתוח $(a,b)$.
  \item $f(x)$ פונקציה רציפה בקטע סגור $[a,b]$.
  \item $f(x)$ פונקציה מונוטונית עולה בקטע פתוח $(a,b)$.
  \item $f(x)$ פונקציה מונוטונית עולה בקטע סגור $[a,b]$.
\end{enumerate}
\end{question}

\begin{hint}
מה אומר משפט לגרנז' על $f(x_2)-f(x_1)$ כאשר $a<x_1<x_2<b$?
\end{hint}

\begin{solution}
\textbf{התשובה הנכונה: ד.}

יהיו $a<x_1<x_2<b$. $f$ גזירה ולכן רציפה ב-$[x_1,x_2]$ וגזירה ב-$(x_1,x_2)$, ולפי משפט לגרנז' קיימת $c\in(x_1,x_2)$ כך ש-
\[ f(x_2)-f(x_1)=f'(c)(x_2-x_1)>0 . \]
לכן $f$ עולה (ממש) ב-$(a,b)$.

דוגמאות נגדיות לשאר הטענות, למשל בקטע $(0,1)$:
\begin{itemize}
  \item $f(x)=-\frac1x$: $f'(x)=\frac1{x^2}>0$, אבל $f$ אינה חסומה ב-$(0,1)$ (א' נופלת), אינה מוגדרת ב-$0$ ובפרט אינה רציפה ב-$[0,1]$ (ג' נופלת) ואינה מונוטונית בקטע הסגור (ה' נופלת, $f$ כלל לא מוגדרת שם).
  \item $f(x)=x-5$: $f'=1>0$ אבל $f<0$ בכל $(0,1)$ (ב' נופלת).
\end{itemize}
\end{solution}
